\section{Related Work}
\label{sec:related_work}

\subsection{Talking Head Generation}
\label{sec: rw_1}
In recent years, audio-driven talking head generation methods have attracted significant attention and provide valuable insight for video dubbing techniques. Talking head generation can be categorized into NeRF-based~\cite{peng2024synctalk, tang2022real, liu2022semantic, guo2021ad}, GAN-based~\cite{chen2019hierarchical, zhou2019talking, zhou2021pose, meshry2021learned, das2020speech}, and diffusion-based~\cite{tian2024emo, chen2024echomimic, xu2025vasa, jiang2024loopy, xu2024hallo, chencafe} approaches.

NeRF-based methods require identity-specific videos for training and additional rendering~\cite{mildenhall2021nerf} time, with early methods like AD-NeRF~\cite{guo2021ad} taking several seconds per frame. Despite the real-time rendering achieved by integrating Instant-NGP~\cite{muller2022instant, tang2022real}, retraining is needed for new identities. Diffusion-based methods, such as EMO Portrait~\cite{tian2024emo}, employ a two-stage training process by integrating ReferenceNet~\cite{hu2024animate}, temporal layers~\cite{guo2023animatediff}, and audio attention layers into Stable Diffusion models~\cite{rombach2022high}. Similar strategies are used in LOOPY~\cite{jiang2024loopy}, Hallo~\cite{xu2024hallo}, and EchoMimic~\cite{chen2024echomimic}. 
These methods allow for one-shot vivid talking head generation from images, but are computationally expensive. 

In contrast, GAN-based methods generate images in one step.
Early methods~\cite{chen2018lip, chen2019hierarchical, zhou2019talking} fail to maintain identity consistency and accurate lip movement. 
To address this, methods like MakeItTalk~\cite{zhou2020makelttalk} and SadTalker~\cite{zhang2023sadtalker} adopt multistage inference, separating audio-to-motion and motion-to-video modeling. 
Although this improves the results, it increases the computational overhead and complexity.

\begin{figure*}[tbp]
  \centering
   \includegraphics[width=\linewidth]{image/fig_stage1.pdf}
   \caption{
   Illustration of MuseTalk's framework. 
   We first encode a reference facial image and an occluded lower half target image into perceptually equivalent latent space. 
   Subsequently, we employ a multimodal U-Net to effectively fuse audio and visual features at various scales. 
   Consequently, the decoded results from the latent space yield more realistic and lip-synced talking face visual content.}
   \label{fig:fig_pipeline}
\end{figure*}

\subsection{Video Dubbing}
\label{sec: rw_2}
Video dubbing focuses on replacing the mouth region of a source face based on driving audio. The most common approach~\cite{prajwal2020lip} involves using a mask in the lower half of the face or the lip region to guide the model to create new visual effects only in these areas, as shown in~\cref{fig:fig_pipeline}. 
Early methods~\cite{prajwal2020lip, zhang2023dinet, cheng2022videoretalking, wang2023seeing, guan2023stylesync, sun2022masked, zhong2023identity} predominantly relied on GANs. 
Although these methods achieved lip movements that were relatively consistent with the audio content, they often struggled with maintaining identity consistency and reconstructing clear teeth details. 
DI-Net~\cite{zhang2023dinet} attempted to train models on high-resolution data, which compromised lip accuracy. 
StyleSync~\cite{guan2023stylesync} highlighted this dilemma, noting that forcing the model to restore too many lip-relevant details might interfere with learning.

Recently, methods such as LatentSync~\cite{li2024latentsync} and DiffTalk~\cite{shen2023difftalk} have leveraged the strong detail generation capabilities of the Latent Diffusion paradigm~\cite{rombach2022high} for video dubbing tasks. 
LatentSync integrates the SyncNet loss~\cite{prajwal2020lip} to enhance lip motion-audio alignment during the one-step denoising process. 
Despite improvements in audio-visual synchornize and clarity, these methods suffer from the heavy inference burden.


